\begin{tabularx}{\textwidth}{@{}>{\raggedright\arraybackslash}p{38mm}>{\raggedright\arraybackslash}X@{}}
\toprule
\textbf{Control} & \textbf{What it does} \\
\midrule
\ann{the map} & The plan the aircraft holds in white, the leg or hold being flown in magenta, the lost-link route in amber dashes, landing sites in green, other aircraft as arrowheads, the camera's target as a green cross, and in blue dashes whatever is being prepared and not yet sent. Drag to move it, the wheel to zoom. A click gives a point to a command that waits for one. \\
\ann{follow} & Keeps the aircraft in the middle of the map. Dragging the map switches it off. \\
\ann{plan} & Shows the whole plan. \\
\ann{north up / track up} & Turns the map: north at the top, or the way the aircraft is going (its track) at the top. The arrow in the corner points to north. \\
\ann{the height profile} & Under the map: the heights the plan asks for, dashed, and the height forecast. \\
\ann{Layers} & Opens and closes the panel of what the map shows: the base map, airspace, weather, airports, link range and the rest. Each layer is Off, Auto or On: Auto shows it when its reason holds and says why, in a chip under the toolbar (the cross on a chip puts that layer away until its reason has gone and come back); On and Off are yours. Escape closes the panel. \\
\ann{Declutter} & Puts every optional layer away; a second click brings them back. The plan, the aircraft and the lost-link route stay. \\
\bottomrule
\end{tabularx}
